% --- Archon physics formalization source begin ---
% archon:physics
% archon:covers PhyXMiniProblems/problem_phyx_mini_0454.lean
% archon:source-report reports/phyx_mini/problem_phyx_mini_0454.source.json

\chapter{Physics problem phyx\_mini\_0454}
\label{ch:PhyXMiniProblems_problem_phyx_mini_0454}

\paragraph{Problem source.}
\#\# Physical scenario

The piston/cylinder in figure contains \$0.1 \textbackslash{}, \textbackslash{}text\{kg\}\$ water at \$500 \textbackslash{}, \textasciicircum{}\{\textbackslash{}circ\}\textbackslash{}text\{C\}\$, \$1000 \textbackslash{}, \textbackslash{}text\{kPa\}\$. The piston has a stop at half of the original volume. The water now cools to a room temperature of \$25 \textbackslash{}, \textasciicircum{}\{\textbackslash{}circ\}\textbackslash{}text\{C\}\$.

\#\# Question

Find the heat transfer in the process.

\#\# Answer choices

- A: A: 0 kJ

- B: B: -114.2 kJ

- C: C: -17.7 kJ

- D: D: -318.72 kJ

\#\# Figure caption (auxiliary; use the image as primary evidence)

The image depicts a cross-sectional diagram of a container or piston system containing water. 

- **Container Structure**: The container is shown with thick walls and a movable piston inside it.
- **Piston**: The piston is labeled with "mₚ", indicating mass or weight, and it rests directly on top of the water.
- **Water**: A labeled section of the container shows it is filled with water. The water occupies the lower portion of the system.
- **Pressure**: Above the piston, there's a label "P₀", signifying external pressure acting on the system.
- **Text Labels**: The diagram includes text annotations "mₚ" for the piston, "Water" for the liquid, and "P₀" for the pressure.

This setup is typical for illustrating principles involving pressure and fluid mechanics.

\#\# Dataset metadata

Temperature and Heat Transfer; Physical Model Grounding Reasoning, Multi-Formula Reasoning

\paragraph{Recorded answer/context.}
D: D: -318.72 kJ

\paragraph{Figure/image path.}
/root/proposal\_for\_physic/science-mango/phyx\_mini\_run/phyx\_data/test\_image/454.png

\paragraph{Formalization target.}
create a compiling Lean file with sorry bodies at `PhyXMiniProblems/problem\_phyx\_mini\_0454.lean`. The Lean declarations must preserve the physical quantities, dimensions or dimensional roles, figure labels, governing-law hypotheses, and final relation expressed by this problem.
Use Mathlib/Physlib names found through LeanExplore where available. If a physics API is missing, introduce faithful local abstractions rather than scalar placeholder aliases.

\begin{theorem}[Physics formalization target]
\label{thm:physics:phyx_mini_0454:target}
\lean{PhyXMiniProblems.ProblemPhyXMini0454.problem_phyx_mini_0454}
\uses{def:physics:phyx-mini-0454:phyxminiproblems-problemphyxmini0454-pistoncylindercoolingsetup, def:physics:phyx-mini-0454:phyxminiproblems-problemphyxmini0454-totalheattransferintowaterinkilojoules, def:physics:phyx-mini-0454:phyxminiproblems-problemphyxmini0454-matchespistoncylindercoolingscenario, def:physics:phyx-mini-0454:phyxminiproblems-problemphyxmini0454-matchessuppliedpistoncylinderfigure, def:physics:phyx-mini-0454:phyxminiproblems-problemphyxmini0454-matchesproblemreadouts, def:physics:phyx-mini-0454:phyxminiproblems-problemphyxmini0454-useswaterpropertytablecalibration, def:physics:phyx-mini-0454:phyxminiproblems-problemphyxmini0454-hasphysicalpistoncylinderparameters, def:physics:phyx-mini-0454:phyxminiproblems-problemphyxmini0454-satisfiesloadedpistonwaterlaws, def:physics:phyx-mini-0454:phyxminiproblems-problemphyxmini0454-displayedheattransferinkilojoules, def:physics:phyx-mini-0454:phyxminiproblems-problemphyxmini0454-recordeddatasetanswer, def:physics:phyx-mini-0454:phyxminiproblems-problemphyxmini0454-roundstodisplayedhundredth, def:physics:phyx-mini-0454:phyxminiproblems-problemphyxmini0454-isuniquematchinganswerchoice}
The assigned autoformalize agent should translate this physics problem into Lean declarations in the covered file, with theorem and lemma proof bodies written as `by sorry`.
\end{theorem}
\begin{proof}
This is an autoformalization task, not a proof task. Produce statements that can later be proved by the physics prover without weakening the physical model.
\end{proof}

% --- Archon declaration topology begin ---
\section{Declaration topology}
\begin{definition}[Mass Quantity]
  \label{def:physics:phyx-mini-0454:phyxminiproblems-problemphyxmini0454-massquantity}
  \lean{PhyXMiniProblems.ProblemPhyXMini0454.MassQuantity}
  A nonnegative physical mass with dimension M
\end{definition}
\begin{proof}
  This is the declared model definition of Mass Quantity.
\end{proof}
\begin{definition}[Volume Quantity]
  \label{def:physics:phyx-mini-0454:phyxminiproblems-problemphyxmini0454-volumequantity}
  \lean{PhyXMiniProblems.ProblemPhyXMini0454.VolumeQuantity}
  A nonnegative physical volume with dimension L³
\end{definition}
\begin{proof}
  This is the declared model definition of Volume Quantity.
\end{proof}
\begin{definition}[Specific Volume Quantity]
  \label{def:physics:phyx-mini-0454:phyxminiproblems-problemphyxmini0454-specificvolumequantity}
  \lean{PhyXMiniProblems.ProblemPhyXMini0454.SpecificVolumeQuantity}
  A nonnegative specific volume with dimension L³ M⁻¹
\end{definition}
\begin{proof}
  This is the declared model definition of Specific Volume Quantity.
\end{proof}
\begin{definition}[Specific Energy Quantity]
  \label{def:physics:phyx-mini-0454:phyxminiproblems-problemphyxmini0454-specificenergyquantity}
  \lean{PhyXMiniProblems.ProblemPhyXMini0454.SpecificEnergyQuantity}
  Signed specific energy with dimension L² T⁻²
\end{definition}
\begin{proof}
  This is the declared model definition of Specific Energy Quantity.
\end{proof}
\begin{definition}[Acceleration Quantity]
  \label{def:physics:phyx-mini-0454:phyxminiproblems-problemphyxmini0454-accelerationquantity}
  \lean{PhyXMiniProblems.ProblemPhyXMini0454.AccelerationQuantity}
  A nonnegative acceleration magnitude with dimension L T⁻²
\end{definition}
\begin{proof}
  This is the declared model definition of Acceleration Quantity.
\end{proof}
\begin{definition}[Measured Temperature]
  \label{def:physics:phyx-mini-0454:phyxminiproblems-problemphyxmini0454-measuredtemperature}
  \lean{PhyXMiniProblems.ProblemPhyXMini0454.MeasuredTemperature}
  An absolute temperature together with the zero-preserving unit in which its magnitude is stored. Celsius is used only as an affine scalar readout
\end{definition}
\begin{proof}
  This is the declared model definition of Measured Temperature.
\end{proof}
\begin{definition}[mass In Kilograms]
  \label{def:physics:phyx-mini-0454:phyxminiproblems-problemphyxmini0454-massinkilograms}
  \lean{PhyXMiniProblems.ProblemPhyXMini0454.massInKilograms}
  \uses{def:physics:phyx-mini-0454:phyxminiproblems-problemphyxmini0454-massquantity}
  Kilogram readout of a physical mass
\end{definition}
\begin{proof}
  This is the declared model definition of mass In Kilograms.
\end{proof}
\begin{definition}[pressure In Pascals]
  \label{def:physics:phyx-mini-0454:phyxminiproblems-problemphyxmini0454-pressureinpascals}
  \lean{PhyXMiniProblems.ProblemPhyXMini0454.pressureInPascals}
  Pascal readout of a physical pressure
\end{definition}
\begin{proof}
  This is the declared model definition of pressure In Pascals.
\end{proof}
\begin{definition}[pressure In Kilopascals]
  \label{def:physics:phyx-mini-0454:phyxminiproblems-problemphyxmini0454-pressureinkilopascals}
  \lean{PhyXMiniProblems.ProblemPhyXMini0454.pressureInKilopascals}
  \uses{def:physics:phyx-mini-0454:phyxminiproblems-problemphyxmini0454-pressureinpascals}
  Kilopascal readout of a physical pressure
\end{definition}
\begin{proof}
  This is the declared model definition of pressure In Kilopascals.
\end{proof}
\begin{definition}[volume In Cubic Metres]
  \label{def:physics:phyx-mini-0454:phyxminiproblems-problemphyxmini0454-volumeincubicmetres}
  \lean{PhyXMiniProblems.ProblemPhyXMini0454.volumeInCubicMetres}
  \uses{def:physics:phyx-mini-0454:phyxminiproblems-problemphyxmini0454-volumequantity}
  Cubic-metre readout of a physical volume
\end{definition}
\begin{proof}
  This is the declared model definition of volume In Cubic Metres.
\end{proof}
\begin{definition}[specific Volume In Cubic Metres Per Kilogram]
  \label{def:physics:phyx-mini-0454:phyxminiproblems-problemphyxmini0454-specificvolumeincubicmetresperkilogram}
  \lean{PhyXMiniProblems.ProblemPhyXMini0454.specificVolumeInCubicMetresPerKilogram}
  \uses{def:physics:phyx-mini-0454:phyxminiproblems-problemphyxmini0454-specificvolumequantity}
  Cubic-metre-per-kilogram readout of a physical specific volume
\end{definition}
\begin{proof}
  This is the declared model definition of specific Volume In Cubic Metres Per Kilogram.
\end{proof}
\begin{definition}[energy In Joules]
  \label{def:physics:phyx-mini-0454:phyxminiproblems-problemphyxmini0454-energyinjoules}
  \lean{PhyXMiniProblems.ProblemPhyXMini0454.energyInJoules}
  Joule readout of physical heat, work, or internal energy
\end{definition}
\begin{proof}
  This is the declared model definition of energy In Joules.
\end{proof}
\begin{definition}[energy In Kilojoules]
  \label{def:physics:phyx-mini-0454:phyxminiproblems-problemphyxmini0454-energyinkilojoules}
  \lean{PhyXMiniProblems.ProblemPhyXMini0454.energyInKilojoules}
  \uses{def:physics:phyx-mini-0454:phyxminiproblems-problemphyxmini0454-energyinjoules}
  Kilojoule readout of physical heat, work, or internal energy
\end{definition}
\begin{proof}
  This is the declared model definition of energy In Kilojoules.
\end{proof}
\begin{definition}[specific Energy In Joules Per Kilogram]
  \label{def:physics:phyx-mini-0454:phyxminiproblems-problemphyxmini0454-specificenergyinjoulesperkilogram}
  \lean{PhyXMiniProblems.ProblemPhyXMini0454.specificEnergyInJoulesPerKilogram}
  \uses{def:physics:phyx-mini-0454:phyxminiproblems-problemphyxmini0454-specificenergyquantity}
  Joule-per-kilogram readout of a physical specific energy
\end{definition}
\begin{proof}
  This is the declared model definition of specific Energy In Joules Per Kilogram.
\end{proof}
\begin{definition}[specific Energy In Kilojoules Per Kilogram]
  \label{def:physics:phyx-mini-0454:phyxminiproblems-problemphyxmini0454-specificenergyinkilojoulesperkilogram}
  \lean{PhyXMiniProblems.ProblemPhyXMini0454.specificEnergyInKilojoulesPerKilogram}
  \uses{def:physics:phyx-mini-0454:phyxminiproblems-problemphyxmini0454-specificenergyquantity, def:physics:phyx-mini-0454:phyxminiproblems-problemphyxmini0454-specificenergyinjoulesperkilogram}
  Kilojoule-per-kilogram readout of a physical specific energy
\end{definition}
\begin{proof}
  This is the declared model definition of specific Energy In Kilojoules Per Kilogram.
\end{proof}
\begin{definition}[area In Square Metres]
  \label{def:physics:phyx-mini-0454:phyxminiproblems-problemphyxmini0454-areainsquaremetres}
  \lean{PhyXMiniProblems.ProblemPhyXMini0454.areaInSquareMetres}
  Square-metre readout of the piston face area
\end{definition}
\begin{proof}
  This is the declared model definition of area In Square Metres.
\end{proof}
\begin{definition}[acceleration In Metres Per Second Squared]
  \label{def:physics:phyx-mini-0454:phyxminiproblems-problemphyxmini0454-accelerationinmetrespersecondsquared}
  \lean{PhyXMiniProblems.ProblemPhyXMini0454.accelerationInMetresPerSecondSquared}
  \uses{def:physics:phyx-mini-0454:phyxminiproblems-problemphyxmini0454-accelerationquantity}
  Metre-per-second-squared readout of an acceleration magnitude
\end{definition}
\begin{proof}
  This is the declared model definition of acceleration In Metres Per Second Squared.
\end{proof}
\begin{definition}[temperature In Kelvin]
  \label{def:physics:phyx-mini-0454:phyxminiproblems-problemphyxmini0454-temperatureinkelvin}
  \lean{PhyXMiniProblems.ProblemPhyXMini0454.temperatureInKelvin}
  \uses{def:physics:phyx-mini-0454:phyxminiproblems-problemphyxmini0454-measuredtemperature}
  Kelvin readout of a measured absolute temperature
\end{definition}
\begin{proof}
  This is the declared model definition of temperature In Kelvin.
\end{proof}
\begin{definition}[temperature In Degrees Celsius]
  \label{def:physics:phyx-mini-0454:phyxminiproblems-problemphyxmini0454-temperatureindegreescelsius}
  \lean{PhyXMiniProblems.ProblemPhyXMini0454.temperatureInDegreesCelsius}
  \uses{def:physics:phyx-mini-0454:phyxminiproblems-problemphyxmini0454-measuredtemperature, def:physics:phyx-mini-0454:phyxminiproblems-problemphyxmini0454-temperatureinkelvin}
  Celsius readout, using the exact affine offset 273.15 K
\end{definition}
\begin{proof}
  This is the declared model definition of temperature In Degrees Celsius.
\end{proof}
\begin{definition}[Process Point]
  \label{def:physics:phyx-mini-0454:phyxminiproblems-problemphyxmini0454-processpoint}
  \lean{PhyXMiniProblems.ProblemPhyXMini0454.ProcessPoint}
  The three equilibrium states needed for the two-stage cooling path
\end{definition}
\begin{proof}
  This is the declared model definition of Process Point.
\end{proof}
\begin{definition}[Process Leg]
  \label{def:physics:phyx-mini-0454:phyxminiproblems-problemphyxmini0454-processleg}
  \lean{PhyXMiniProblems.ProblemPhyXMini0454.ProcessLeg}
  The two consecutive legs of the cooling process
\end{definition}
\begin{proof}
  This is the declared model definition of Process Leg.
\end{proof}
\begin{definition}[leg Start]
  \label{def:physics:phyx-mini-0454:phyxminiproblems-problemphyxmini0454-legstart}
  \lean{PhyXMiniProblems.ProblemPhyXMini0454.legStart}
  \uses{def:physics:phyx-mini-0454:phyxminiproblems-problemphyxmini0454-processpoint, def:physics:phyx-mini-0454:phyxminiproblems-problemphyxmini0454-processleg}
  Initial point of each directed process leg
\end{definition}
\begin{proof}
  This is the declared model definition of leg Start.
\end{proof}
\begin{definition}[leg Finish]
  \label{def:physics:phyx-mini-0454:phyxminiproblems-problemphyxmini0454-legfinish}
  \lean{PhyXMiniProblems.ProblemPhyXMini0454.legFinish}
  \uses{def:physics:phyx-mini-0454:phyxminiproblems-problemphyxmini0454-processpoint, def:physics:phyx-mini-0454:phyxminiproblems-problemphyxmini0454-processleg}
  Final point of each directed process leg
\end{definition}
\begin{proof}
  This is the declared model definition of leg Finish.
\end{proof}
\begin{definition}[Process Constraint]
  \label{def:physics:phyx-mini-0454:phyxminiproblems-problemphyxmini0454-processconstraint}
  \lean{PhyXMiniProblems.ProblemPhyXMini0454.ProcessConstraint}
  Mechanical constraint on a leg of the process
\end{definition}
\begin{proof}
  This is the declared model definition of Process Constraint.
\end{proof}
\begin{definition}[expected Constraint]
  \label{def:physics:phyx-mini-0454:phyxminiproblems-problemphyxmini0454-expectedconstraint}
  \lean{PhyXMiniProblems.ProblemPhyXMini0454.expectedConstraint}
  \uses{def:physics:phyx-mini-0454:phyxminiproblems-problemphyxmini0454-processleg, def:physics:phyx-mini-0454:phyxminiproblems-problemphyxmini0454-processconstraint}
  Constraint dictated by the loaded piston and the rigid stops
\end{definition}
\begin{proof}
  This is the declared model definition of expected Constraint.
\end{proof}
\begin{definition}[Working Fluid]
  \label{def:physics:phyx-mini-0454:phyxminiproblems-problemphyxmini0454-workingfluid}
  \lean{PhyXMiniProblems.ProblemPhyXMini0454.WorkingFluid}
  The working substance named in the problem and raster
\end{definition}
\begin{proof}
  This is the declared model definition of Working Fluid.
\end{proof}
\begin{definition}[Figure Label]
  \label{def:physics:phyx-mini-0454:phyxminiproblems-problemphyxmini0454-figurelabel}
  \lean{PhyXMiniProblems.ProblemPhyXMini0454.FigureLabel}
  Literal symbolic or text labels visible in the primary raster
\end{definition}
\begin{proof}
  This is the declared model definition of Figure Label.
\end{proof}
\begin{definition}[Piston Cylinder Figure]
  \label{def:physics:phyx-mini-0454:phyxminiproblems-problemphyxmini0454-pistoncylinderfigure}
  \lean{PhyXMiniProblems.ProblemPhyXMini0454.PistonCylinderFigure}
  \uses{def:physics:phyx-mini-0454:phyxminiproblems-problemphyxmini0454-figurelabel}
  Qualitative geometry and labels read from the supplied image. The numerical half-volume placement comes from the prose and is recorded separately below
\end{definition}
\begin{proof}
  This is the declared model definition of Piston Cylinder Figure.
\end{proof}
\begin{definition}[Water Equilibrium State]
  \label{def:physics:phyx-mini-0454:phyxminiproblems-problemphyxmini0454-waterequilibriumstate}
  \lean{PhyXMiniProblems.ProblemPhyXMini0454.WaterEquilibriumState}
  \uses{def:physics:phyx-mini-0454:phyxminiproblems-problemphyxmini0454-volumequantity, def:physics:phyx-mini-0454:phyxminiproblems-problemphyxmini0454-specificvolumequantity, def:physics:phyx-mini-0454:phyxminiproblems-problemphyxmini0454-specificenergyquantity, def:physics:phyx-mini-0454:phyxminiproblems-problemphyxmini0454-measuredtemperature}
  A physical equilibrium state of the fixed water sample
\end{definition}
\begin{proof}
  This is the declared model definition of Water Equilibrium State.
\end{proof}
\begin{definition}[Piston Cylinder Cooling Setup]
  \label{def:physics:phyx-mini-0454:phyxminiproblems-problemphyxmini0454-pistoncylindercoolingsetup}
  \lean{PhyXMiniProblems.ProblemPhyXMini0454.PistonCylinderCoolingSetup}
  \uses{def:physics:phyx-mini-0454:phyxminiproblems-problemphyxmini0454-massquantity, def:physics:phyx-mini-0454:phyxminiproblems-problemphyxmini0454-specificvolumequantity, def:physics:phyx-mini-0454:phyxminiproblems-problemphyxmini0454-specificenergyquantity, def:physics:phyx-mini-0454:phyxminiproblems-problemphyxmini0454-accelerationquantity, def:physics:phyx-mini-0454:phyxminiproblems-problemphyxmini0454-measuredtemperature, def:physics:phyx-mini-0454:phyxminiproblems-problemphyxmini0454-processpoint, def:physics:phyx-mini-0454:phyxminiproblems-problemphyxmini0454-processleg, def:physics:phyx-mini-0454:phyxminiproblems-problemphyxmini0454-processconstraint, def:physics:phyx-mini-0454:phyxminiproblems-problemphyxmini0454-workingfluid, def:physics:phyx-mini-0454:phyxminiproblems-problemphyxmini0454-pistoncylinderfigure, def:physics:phyx-mini-0454:phyxminiproblems-problemphyxmini0454-waterequilibriumstate}
  Independent observables and material-property functions for the apparatus. Neither total heat transfer nor a selected answer is stored as a field
\end{definition}
\begin{proof}
  This is the declared model definition of Piston Cylinder Cooling Setup.
\end{proof}
\begin{definition}[loaded Piston Pressure In Pascals]
  \label{def:physics:phyx-mini-0454:phyxminiproblems-problemphyxmini0454-loadedpistonpressureinpascals}
  \lean{PhyXMiniProblems.ProblemPhyXMini0454.loadedPistonPressureInPascals}
  \uses{def:physics:phyx-mini-0454:phyxminiproblems-problemphyxmini0454-massinkilograms, def:physics:phyx-mini-0454:phyxminiproblems-problemphyxmini0454-pressureinpascals, def:physics:phyx-mini-0454:phyxminiproblems-problemphyxmini0454-areainsquaremetres, def:physics:phyx-mini-0454:phyxminiproblems-problemphyxmini0454-accelerationinmetrespersecondsquared, def:physics:phyx-mini-0454:phyxminiproblems-problemphyxmini0454-pistoncylindercoolingsetup}
  Pressure due to P₀ together with the piston weight, in pascals
\end{definition}
\begin{proof}
  This is the declared model definition of loaded Piston Pressure In Pascals.
\end{proof}
\begin{definition}[total Boundary Work In Kilojoules]
  \label{def:physics:phyx-mini-0454:phyxminiproblems-problemphyxmini0454-totalboundaryworkinkilojoules}
  \lean{PhyXMiniProblems.ProblemPhyXMini0454.totalBoundaryWorkInKilojoules}
  \uses{def:physics:phyx-mini-0454:phyxminiproblems-problemphyxmini0454-energyinkilojoules, def:physics:phyx-mini-0454:phyxminiproblems-problemphyxmini0454-processleg, def:physics:phyx-mini-0454:phyxminiproblems-problemphyxmini0454-pistoncylindercoolingsetup}
  Total boundary work done by the water during both legs, in kilojoules
\end{definition}
\begin{proof}
  This is the declared model definition of total Boundary Work In Kilojoules.
\end{proof}
\begin{definition}[total Heat Transfer Into Water In Kilojoules]
  \label{def:physics:phyx-mini-0454:phyxminiproblems-problemphyxmini0454-totalheattransferintowaterinkilojoules}
  \lean{PhyXMiniProblems.ProblemPhyXMini0454.totalHeatTransferIntoWaterInKilojoules}
  \uses{def:physics:phyx-mini-0454:phyxminiproblems-problemphyxmini0454-energyinkilojoules, def:physics:phyx-mini-0454:phyxminiproblems-problemphyxmini0454-processleg, def:physics:phyx-mini-0454:phyxminiproblems-problemphyxmini0454-pistoncylindercoolingsetup}
  Total signed heat transferred into the water during both legs, in kJ
\end{definition}
\begin{proof}
  This is the declared model definition of total Heat Transfer Into Water In Kilojoules.
\end{proof}
\begin{definition}[internal Energy Change In Kilojoules]
  \label{def:physics:phyx-mini-0454:phyxminiproblems-problemphyxmini0454-internalenergychangeinkilojoules}
  \lean{PhyXMiniProblems.ProblemPhyXMini0454.internalEnergyChangeInKilojoules}
  \uses{def:physics:phyx-mini-0454:phyxminiproblems-problemphyxmini0454-energyinkilojoules, def:physics:phyx-mini-0454:phyxminiproblems-problemphyxmini0454-pistoncylindercoolingsetup}
  Total internal-energy change of the water, in kilojoules
\end{definition}
\begin{proof}
  This is the declared model definition of internal Energy Change In Kilojoules.
\end{proof}
\begin{definition}[Matches Piston Cylinder Cooling Scenario]
  \label{def:physics:phyx-mini-0454:phyxminiproblems-problemphyxmini0454-matchespistoncylindercoolingscenario}
  \lean{PhyXMiniProblems.ProblemPhyXMini0454.MatchesPistonCylinderCoolingScenario}
  \uses{def:physics:phyx-mini-0454:phyxminiproblems-problemphyxmini0454-expectedconstraint, def:physics:phyx-mini-0454:phyxminiproblems-problemphyxmini0454-pistoncylindercoolingsetup}
  Qualitative closed-system and two-stage process assumptions
\end{definition}
\begin{proof}
  This is the declared model definition of Matches Piston Cylinder Cooling Scenario.
\end{proof}
\begin{definition}[Matches Supplied Piston Cylinder Figure]
  \label{def:physics:phyx-mini-0454:phyxminiproblems-problemphyxmini0454-matchessuppliedpistoncylinderfigure}
  \lean{PhyXMiniProblems.ProblemPhyXMini0454.MatchesSuppliedPistonCylinderFigure}
  \uses{def:physics:phyx-mini-0454:phyxminiproblems-problemphyxmini0454-pistoncylindercoolingsetup}
  Primary-image evidence: the raster shows a water-filled cylinder, a movable piston labelled mₚ, an external pressure P₀, and opposed internal stops
\end{definition}
\begin{proof}
  This is the declared model definition of Matches Supplied Piston Cylinder Figure.
\end{proof}
\begin{definition}[Matches Problem Readouts]
  \label{def:physics:phyx-mini-0454:phyxminiproblems-problemphyxmini0454-matchesproblemreadouts}
  \lean{PhyXMiniProblems.ProblemPhyXMini0454.MatchesProblemReadouts}
  \uses{def:physics:phyx-mini-0454:phyxminiproblems-problemphyxmini0454-massinkilograms, def:physics:phyx-mini-0454:phyxminiproblems-problemphyxmini0454-pressureinkilopascals, def:physics:phyx-mini-0454:phyxminiproblems-problemphyxmini0454-volumeincubicmetres, def:physics:phyx-mini-0454:phyxminiproblems-problemphyxmini0454-temperatureindegreescelsius, def:physics:phyx-mini-0454:phyxminiproblems-problemphyxmini0454-pistoncylindercoolingsetup}
  Numerical data stated in the prose. The stop and final states have half the initial total volume; no heat-transfer value occurs in these premises
\end{definition}
\begin{proof}
  This is the declared model definition of Matches Problem Readouts.
\end{proof}
\begin{definition}[Uses Water Property Table Calibration]
  \label{def:physics:phyx-mini-0454:phyxminiproblems-problemphyxmini0454-useswaterpropertytablecalibration}
  \lean{PhyXMiniProblems.ProblemPhyXMini0454.UsesWaterPropertyTableCalibration}
  \uses{def:physics:phyx-mini-0454:phyxminiproblems-problemphyxmini0454-pressureinkilopascals, def:physics:phyx-mini-0454:phyxminiproblems-problemphyxmini0454-specificvolumeincubicmetresperkilogram, def:physics:phyx-mini-0454:phyxminiproblems-problemphyxmini0454-specificenergyinkilojoulesperkilogram, def:physics:phyx-mini-0454:phyxminiproblems-problemphyxmini0454-pistoncylindercoolingsetup}
  Rounded superheated-water and saturated-water property-table entries used in the calculation. These are independent material data, not a heat-transfer answer. At 1000 kPa, 500 °C, v₁ = 0.3541 m³/kg and u₁ = 3124.3 kJ/kg. At 25 °C, the listed saturated endpoints determine the final mixture quality and energy through the constitutive laws below
\end{definition}
\begin{proof}
  This is the declared model definition of Uses Water Property Table Calibration.
\end{proof}
\begin{definition}[Has Physical Piston Cylinder Parameters]
  \label{def:physics:phyx-mini-0454:phyxminiproblems-problemphyxmini0454-hasphysicalpistoncylinderparameters}
  \lean{PhyXMiniProblems.ProblemPhyXMini0454.HasPhysicalPistonCylinderParameters}
  \uses{def:physics:phyx-mini-0454:phyxminiproblems-problemphyxmini0454-massinkilograms, def:physics:phyx-mini-0454:phyxminiproblems-problemphyxmini0454-pressureinpascals, def:physics:phyx-mini-0454:phyxminiproblems-problemphyxmini0454-volumeincubicmetres, def:physics:phyx-mini-0454:phyxminiproblems-problemphyxmini0454-areainsquaremetres, def:physics:phyx-mini-0454:phyxminiproblems-problemphyxmini0454-accelerationinmetrespersecondsquared, def:physics:phyx-mini-0454:phyxminiproblems-problemphyxmini0454-pistoncylindercoolingsetup}
  Positivity and phase-range conditions selecting a physical process
\end{definition}
\begin{proof}
  This is the declared model definition of Has Physical Piston Cylinder Parameters.
\end{proof}
\begin{definition}[Satisfies Loaded Piston Water Laws]
  \label{def:physics:phyx-mini-0454:phyxminiproblems-problemphyxmini0454-satisfiesloadedpistonwaterlaws}
  \lean{PhyXMiniProblems.ProblemPhyXMini0454.SatisfiesLoadedPistonWaterLaws}
  \uses{def:physics:phyx-mini-0454:phyxminiproblems-problemphyxmini0454-massinkilograms, def:physics:phyx-mini-0454:phyxminiproblems-problemphyxmini0454-pressureinpascals, def:physics:phyx-mini-0454:phyxminiproblems-problemphyxmini0454-volumeincubicmetres, def:physics:phyx-mini-0454:phyxminiproblems-problemphyxmini0454-specificvolumeincubicmetresperkilogram, def:physics:phyx-mini-0454:phyxminiproblems-problemphyxmini0454-energyinjoules, def:physics:phyx-mini-0454:phyxminiproblems-problemphyxmini0454-specificenergyinjoulesperkilogram, def:physics:phyx-mini-0454:phyxminiproblems-problemphyxmini0454-specificenergyinkilojoulesperkilogram, def:physics:phyx-mini-0454:phyxminiproblems-problemphyxmini0454-legstart, def:physics:phyx-mini-0454:phyxminiproblems-problemphyxmini0454-legfinish, def:physics:phyx-mini-0454:phyxminiproblems-problemphyxmini0454-pistoncylindercoolingsetup, def:physics:phyx-mini-0454:phyxminiproblems-problemphyxmini0454-loadedpistonpressureinpascals}
  Governing mechanics and thermodynamics: total volume and internal energy are mass times their specific values; the freely moving loaded piston maintains the pressure generated by P₀ and mₚ g/A until it reaches the stop; the second leg is isochoric; the final state is a saturated water mixture whose v and u interpolate between saturated liquid and vapor using one vapor quality; quasistatic boundary work is p ΔV; on each closed-system leg, Q\_into = ΔU + W\_by.
\end{definition}
\begin{proof}
  This is the declared model definition of Satisfies Loaded Piston Water Laws.
\end{proof}
\begin{definition}[Answer Choice]
  \label{def:physics:phyx-mini-0454:phyxminiproblems-problemphyxmini0454-answerchoice}
  \lean{PhyXMiniProblems.ProblemPhyXMini0454.AnswerChoice}
  Labels of the four heat-transfer choices printed with the problem
\end{definition}
\begin{proof}
  This is the declared model definition of Answer Choice.
\end{proof}
\begin{definition}[displayed Heat Transfer In Kilojoules]
  \label{def:physics:phyx-mini-0454:phyxminiproblems-problemphyxmini0454-displayedheattransferinkilojoules}
  \lean{PhyXMiniProblems.ProblemPhyXMini0454.displayedHeatTransferInKilojoules}
  \uses{def:physics:phyx-mini-0454:phyxminiproblems-problemphyxmini0454-answerchoice}
  Signed heat transfer in kilojoules printed beside each answer label
\end{definition}
\begin{proof}
  This is the declared model definition of displayed Heat Transfer In Kilojoules.
\end{proof}
\begin{definition}[recorded Dataset Answer]
  \label{def:physics:phyx-mini-0454:phyxminiproblems-problemphyxmini0454-recordeddatasetanswer}
  \lean{PhyXMiniProblems.ProblemPhyXMini0454.recordedDatasetAnswer}
  \uses{def:physics:phyx-mini-0454:phyxminiproblems-problemphyxmini0454-answerchoice}
  The source dataset records choice D; this metadata is not a premise
\end{definition}
\begin{proof}
  This is the declared model definition of recorded Dataset Answer.
\end{proof}
\begin{definition}[Rounds To Displayed Hundredth]
  \label{def:physics:phyx-mini-0454:phyxminiproblems-problemphyxmini0454-roundstodisplayedhundredth}
  \lean{PhyXMiniProblems.ProblemPhyXMini0454.RoundsToDisplayedHundredth}
  The actual heat transfer rounds to a displayed value at two decimals
\end{definition}
\begin{proof}
  This is the declared model definition of Rounds To Displayed Hundredth.
\end{proof}
\begin{definition}[Matches Answer Choice]
  \label{def:physics:phyx-mini-0454:phyxminiproblems-problemphyxmini0454-matchesanswerchoice}
  \lean{PhyXMiniProblems.ProblemPhyXMini0454.MatchesAnswerChoice}
  \uses{def:physics:phyx-mini-0454:phyxminiproblems-problemphyxmini0454-pistoncylindercoolingsetup, def:physics:phyx-mini-0454:phyxminiproblems-problemphyxmini0454-totalheattransferintowaterinkilojoules, def:physics:phyx-mini-0454:phyxminiproblems-problemphyxmini0454-answerchoice, def:physics:phyx-mini-0454:phyxminiproblems-problemphyxmini0454-displayedheattransferinkilojoules, def:physics:phyx-mini-0454:phyxminiproblems-problemphyxmini0454-roundstodisplayedhundredth}
  A choice displays the computed total heat transfer to two decimals
\end{definition}
\begin{proof}
  This is the declared model definition of Matches Answer Choice.
\end{proof}
\begin{definition}[Is Unique Matching Answer Choice]
  \label{def:physics:phyx-mini-0454:phyxminiproblems-problemphyxmini0454-isuniquematchinganswerchoice}
  \lean{PhyXMiniProblems.ProblemPhyXMini0454.IsUniqueMatchingAnswerChoice}
  \uses{def:physics:phyx-mini-0454:phyxminiproblems-problemphyxmini0454-pistoncylindercoolingsetup, def:physics:phyx-mini-0454:phyxminiproblems-problemphyxmini0454-answerchoice, def:physics:phyx-mini-0454:phyxminiproblems-problemphyxmini0454-matchesanswerchoice}
  A choice is the unique displayed value matching the computed heat
\end{definition}
\begin{proof}
  This is the declared model definition of Is Unique Matching Answer Choice.
\end{proof}
\begin{lemma}[cooling Process Energy Accounting]
  \label{lem:physics:phyx-mini-0454:phyxminiproblems-problemphyxmini0454-coolingprocessenergyaccounting}
  \lean{PhyXMiniProblems.ProblemPhyXMini0454.coolingProcessEnergyAccounting}
  \uses{def:physics:phyx-mini-0454:phyxminiproblems-problemphyxmini0454-pistoncylindercoolingsetup, def:physics:phyx-mini-0454:phyxminiproblems-problemphyxmini0454-totalboundaryworkinkilojoules, def:physics:phyx-mini-0454:phyxminiproblems-problemphyxmini0454-totalheattransferintowaterinkilojoules, def:physics:phyx-mini-0454:phyxminiproblems-problemphyxmini0454-internalenergychangeinkilojoules, def:physics:phyx-mini-0454:phyxminiproblems-problemphyxmini0454-matchespistoncylindercoolingscenario, def:physics:phyx-mini-0454:phyxminiproblems-problemphyxmini0454-matchesproblemreadouts, def:physics:phyx-mini-0454:phyxminiproblems-problemphyxmini0454-useswaterpropertytablecalibration, def:physics:phyx-mini-0454:phyxminiproblems-problemphyxmini0454-satisfiesloadedpistonwaterlaws, def:physics:phyx-mini-0454:phyxminiproblems-problemphyxmini0454-roundstodisplayedhundredth}
  The moving-piston leg performs about -17.705 kJ of boundary work and the fixed-volume leg performs none. Combining the property-table mixture calculation with the first law gives a total heat transfer that rounds to -318.72 kJ. Both statements are derived conclusions
\end{lemma}
\begin{proof}
  The conclusion follows from the governing relations and figure readouts named in the statement of cooling Process Energy Accounting.
\end{proof}
% --- Archon declaration topology end ---
% --- Archon physics formalization source end ---
